\section*{Capítulo \ref{ch:b1:cross-sections} --- Cortes geológicos}

\begin{solution}{exo:b1:cross-sections:1}
A 1:25\,000, \qty{1}{cm} del mapa son \qty{250}{m} sobre el terreno. La escala vertical es de
\qty{1}{cm} por cada \qty{100}{m}, así que la escala vertical es
$250/100 = 2.5$ veces la horizontal: $e = 2.5$.
\end{solution}

\begin{solution}{exo:b1:cross-sections:2}
$\gamma = \ang{90} - \ang{60} = \ang{30}$, y $\tan\beta = \tan\ang{40}
\cos\ang{30} = 0.8391\times0.8660 = 0.7267$: $\beta = \ang{36.0}$, el estrato
baja hacia el extremo este del corte. En el corte de \omterm{def:b1:spherical-astronomy:horizontal}{acimut}
\ang{000}, $\gamma = \ang{90}$ y $\beta = 0$: el corte sigue el \omterm{def:b1:geological-maps:attitude}{rumbo} y el estrato
aparece horizontal.
\end{solution}

\begin{solution}{exo:b1:cross-sections:3}
$t = h\cos\delta = 120\times\cos\ang{12.16} = 120\times0.9776 =
\qty{117.3}{m}$.
\end{solution}

\begin{solution}{exo:b1:cross-sections:4}
$\tan\ang{12} = 0.2126$. Con $e = 2$: $\arctan 0.4251 = \ang{23.0}$; con
$e = 3$: $\arctan 0.6377 = \ang{32.5}$; con $e = 5$: $\arctan 1.0628 =
\ang{46.7}$. El ángulo crece mucho menos que $e$ y nunca llega a \ang{90}.
\end{solution}

\begin{solution}{exo:b1:cross-sections:5}
(a) $t = 300\sin\ang{25} = \qty{126.8}{m}$. (b) $t = 300\sin\ang{35} =
\qty{172.1}{m}$. (c) $t = 300\sin\ang{15} = \qty{77.6}{m}$. Una misma anchura de afloramiento
puede corresponder a espesores que difieren más de un factor dos: hay que tener siempre en cuenta
la pendiente del terreno.
\end{solution}

\begin{solution}{exo:b1:cross-sections:6}
Con \omterm{def:b1:geological-maps:attitude}{dirección de buzamiento} $\alpha$ y \omterm{def:b1:geological-maps:attitude}{buzamiento} $\delta$, el teorema da en el corte norte--sur
($\theta = 0$) $\tan\ang{20} = \tan\delta\cos\alpha$, y en el corte este--oeste
($\theta = \ang{90}$) $\tan\ang{30} =
\tan\delta\cos(\alpha - \ang{90}) = \tan\delta\sin\alpha$. Dividiendo,
$\tan\alpha = 0.5774/0.3640 = 1.586$, $\alpha = \ang{057.8}$; sumando los cuadrados,
$\tan\delta = \sqrt{0.3640^2 + 0.5774^2} = 0.6825$, $\delta =
\ang{34.3}$. Los dos \omterm{def:b1:cross-sections:apparent-dip}{buzamientos aparentes} son menores que el verdadero, como debe ser.
\end{solution}

\begin{solution}{exo:b1:cross-sections:7}
$100/\tan\ang{12.16} = 100/0.2155 = \qty{464.1}{m}$. En el mapa, las \omterm{def:b1:cross-sections:structure-contour}{curvas de nivel
estructurales} del flanco oriental están, en efecto, a un poco menos de medio kilómetro, equidistantes
y paralelas: el flanco es un plano.
\end{solution}

\begin{solution}{exo:b1:cross-sections:8}
B1: $524.02 - 488.98 = \qty{35.04}{m}$; B3: $529.98 - 494.93 =
\qty{35.05}{m}$. Sobre B2, la base de la lutita es el techo de la caliza, \qty{180}{m} por encima de
su base: $280.16 + 180 = \qty{460.16}{m}$, más alta que el terreno (\qty{380.64}{m}), por eso está
erosionada allí. Entre B1 y B2 el contacto sube $460.16 - 35.04 = \qty{425.12}{m}$ a lo largo de
\qty{2000}{m}: una pendiente de $0.2126$, un \omterm{def:b1:cross-sections:apparent-dip}{buzamiento aparente} de \ang{12.0}. B1 y B3, a la misma
distancia a cada lado de B2, encuentran cada contacto a la misma altitud con \qty{1}{cm} de
diferencia: el \omterm{def:b1:geological-maps:fold}{pliegue} es simétrico respecto de un plano axial vertical en $x = \qty{3.5}{km}$, sus
dos flancos buzan igual, y ninguna falla corta el corte entre ellos.
\end{solution}

\begin{solution}{exo:b1:cross-sections:9}
\omterm{def:b1:cross-sections:throw}{Salto horizontal} $= 150/\tan\ang{50} = 150/1.1918 = \qty{125.9}{m}$. En el corte de \omterm{def:b1:spherical-astronomy:horizontal}{acimut}
\ang{150}, $\gamma = \ang{30}$ y $\tan\beta = 1.1918\times0.8660 =
1.0321$: la falla aparece con el \omterm{def:b1:cross-sections:apparent-dip}{buzamiento aparente} \ang{45.9}. Los cortes de un estrato horizontal
están separados \qty{150}{m} en vertical en esa traza, luego $150/1.0321 = \qty{145.3}{m}$ de
separación a lo largo del corte, más que el \omterm{def:b1:cross-sections:throw}{salto horizontal}, ya que el corte cruza la falla
oblicuamente.
\end{solution}

\begin{solution}{exo:b1:cross-sections:10}
$\cos\gamma = \tan\ang{20}/\tan\ang{40} = 0.3640/0.8391 = 0.4338$, así que
$\gamma = \ang{64.3}$: los cortes de \omterm{def:b1:spherical-astronomy:horizontal}{acimutes} $\alpha \pm \ang{64.3}$ (y las direcciones
opuestas, donde el estrato sube). El \omterm{def:b1:cross-sections:apparent-dip}{buzamiento aparente} es \ang{40.0}, \ang{39.6}, \ang{38.3} y
\ang{36.0} para $\gamma = 0$, \ang{10}, \ang{20}, \ang{30}, pero \ang{22.8}, \ang{16.0}, \ang{8.3} y
$0$ para $\gamma = \ang{60}$, \ang{70}, \ang{80}, \ang{90}. En efecto, derivando
$\tan\beta = \tan\delta\cos\gamma$,
$\dfrac{\mathrm d\beta}{\mathrm d\gamma} =
-\dfrac{\tan\delta\sin\gamma}{1 + \tan^2\delta\cos^2\gamma}$, que es cero
en $\gamma = 0$ y vale $-\tan\delta$ en $\gamma = \ang{90}$: un corte algo apartado de la \omterm{def:b1:geological-maps:attitude}{dirección
de buzamiento} muestra casi el \omterm{def:b1:geological-maps:attitude}{buzamiento} verdadero, mientras que cerca del \omterm{def:b1:geological-maps:attitude}{rumbo} cada grado de \omterm{def:b1:spherical-astronomy:horizontal}{acimut}
cambia el \omterm{def:b1:cross-sections:apparent-dip}{buzamiento aparente} casi un grado.
\end{solution}

\begin{solution}{exo:b1:cross-sections:11}
Los estratos buzan hacia el este, así que la normal a la estratificación, seguida hacia abajo, está
inclinada \ang{12.16} desde la vertical hacia el oeste, hacia \ang{260.7}. (a) Un sondeo vertical
forma con ella el ángulo $\psi = \ang{12.16}$: $176.0/\cos\ang{12.16} = \qty{180.0}{m}$, el \omterm{def:b1:cross-sections:thickness}{espesor
vertical}. (b) Un sondeo que se hunde \ang{60} hacia \ang{260.7} está inclinado \ang{30} desde la
vertical, del mismo lado que la normal: $\psi = \ang{30} - \ang{12.16} =
\ang{17.84}$, $\cos\psi = 0.9519$, y el testigo mide $176.0/0.9519 =
\qty{184.9}{m}$ de longitud. (c) Hacia \ang{080.7} está inclinado al otro lado:
$\psi = \ang{30} + \ang{12.16} = \ang{42.16}$, $\cos\psi = 0.7413$, y el testigo mide
$176.0/0.7413 = \qty{237.4}{m}$. El testigo es más corto, igual a
$t$, cuando el sondeo sigue la normal, hundiéndose \ang{77.84} hacia \ang{260.7}.
\end{solution}

\begin{solution}{exo:b1:cross-sections:12}
En el plano vertical este--oeste, con los ángulos medidos desde el este, los estratos superiores
suben hacia el este con \ang{20} y los inferiores bajan hacia el este con \ang{70}. Girando el
afloramiento alrededor del \omterm{def:b1:geological-maps:attitude}{rumbo} \ang{20}, con el lado este hacia abajo, los estratos superiores
quedan horizontales y los inferiores pasan a $\ang{70} + \ang{20} = \ang{90}$: estaban verticales
cuando se depositaron los superiores. Historia, desde el acontecimiento más antiguo: los estratos
inferiores se depositan planos; se levantan de canto, hasta la vertical, y se elevan; sus bordes se
erosionan hasta una superficie; los estratos superiores se depositan sobre esa superficie; el
conjunto se inclina \ang{20} hacia el oeste; la erosión corta la superficie del terreno de hoy. Es
la historia del afloramiento de Hutton en Siccar Point.
\end{solution}

\begin{solution}{pb:b1:cross-sections:1}
\textbf{1.} El espesor sobre el túnel es $z - 200$: mayor en $x = \qty{1}{km}$,
$556 - 200 = \qty{356}{m}$; menor en $x =
\qty{9}{km}$, $339 - 200 = \qty{139}{m}$.

\textbf{2.} \qty{1}{km} de altitud mide \qty{4.2}{cm}, y \qty{1}{km} de distancia mide
\qty{1.4}{cm}: $e = 4.2/1.4 = 3$.

\textbf{3.} $380.64 - 220.48 = \qty{160.16}{m}$.

\textbf{4.} En B2 la base de la caliza está en $380.64 - 100.48 =
\qty{280.16}{m}$ y la base de la arenisca en \qty{160.16}{m}: el túnel, a \qty{200}{m}, está en la
arenisca.

\textbf{5.} B1: $524.02 - 668.98 = \qty{-144.96}{m}$; B3: $529.98 - 674.93 =
\qty{-144.95}{m}$.

\textbf{6.} De oeste a este: creta, lutita, caliza, arenisca (en el núcleo del \omterm{def:b1:geological-maps:fold}{anticlinal}), caliza,
lutita, y de nuevo creta hasta $x = \qty{9}{km}$, a través del \omterm{def:b1:geological-maps:fold}{sinclinal} y cruzando la falla.

\textbf{7.} $\gamma = \ang{80.7} - \ang{90} = \ang{-9.3}$; $\tan\beta =
\tan\ang{12.16}\cos\ang{9.3} = 0.2155\times0.9869 = 0.2126$: $\beta =
\ang{12.0}$, hacia abajo, hacia el este.

\textbf{8.} $(280.16 + 144.95)/2000 = 425.11/2000 = 0.2126$: \ang{12.0}.

\textbf{9.} $\arctan(3\times0.2126) = \arctan 0.6377 = \ang{32.5}$.

\textbf{10.} $\gamma = \ang{120} - \ang{90} = \ang{30}$; $\tan\beta =
\tan\ang{60}\cos\ang{30} = 1.732\times0.866 = 1.500$: $\beta = \ang{56.3}$.

\textbf{11.} Los flancos buzan hacia \ang{080.7} o \ang{279.3}, solo \ang{9.3} respecto de la línea
del corte, y $\cos\ang{9.3} = 0.987$; la falla buza hacia \ang{120}, \ang{30} respecto de él, y
$\cos\ang{30} = 0.866$.

\textbf{12.} El plano de falla es $z = 250 - \tan\ang{60}\,[(x - 4300)
\sin\ang{120} + (y - 1000)\cos\ang{120}]$. Con $y = 6000$, $(y - 1000)
\cos\ang{120} = -2500$, así que $(x - 4300)\sin\ang{120} = 2500 - (z - 250)/
\tan\ang{60}$, y $1/\tan\ang{60} = \tan\ang{30}$.

\textbf{13.} A $z = 200$: $x = 4300 + (2500 + 50\times0.5774)/0.8660 =
4300 + 2528.9/0.8660 = \qty{7220}{m}$, es decir, \qty{7.22}{km} de C.

\textbf{14.} \omterm{def:b1:cross-sections:throw}{Salto horizontal} $= 200/\tan\ang{60} = \qty{115.5}{m}$; separación a lo largo de
C--D $= 200/1.500 = \qty{133.3}{m}$.

\textbf{15.} La falla buza hacia el estesureste y es normal: el bloque del techo, al este, está
hundido. A la altura del túnel la creta queda a ambos lados de la falla: al oeste, la creta del
\omterm{def:b1:geological-maps:fold}{sinclinal}; al este, creta llevada \qty{200}{m} hacia abajo.

\textbf{16.} El túnel sigue en la creta al cruzar la falla, así que ningún cambio de roca en el frente
lo anuncia; solo lo hace una zona de roca triturada, y quizá agua. Los ingenieros la predicen a partir
del corte (aquí \qty{7.22}{km}), perforan sondeos de reconocimiento por delante de la máquina al
acercarse y preparan el sostenimiento de la zona triturada.

\textbf{17.} $668.98 - 488.98 = \qty{180.00}{m}$.

\textbf{18.} $t = 180\cos\ang{12.16} = \qty{176.0}{m}$.

\textbf{19.} En la cresta ($x = 3500$) la base de la arenisca está en \qty{160.16}{m} (pregunta 3);
baja en ambos flancos con el \omterm{def:b1:cross-sections:apparent-dip}{buzamiento aparente} \ang{12.0} de la pregunta 7, es decir, en
$\tan\ang{12.0}$ por metro de distancia horizontal desde el eje.

\textbf{20.} La base de la caliza está \qty{120}{m} por encima de la base de la arenisca, en
$280.16 - u\tan\ang{12.0}$ con $u = |x - 3500|$, y su techo \qty{180}{m} más arriba, en
$460.16 - u\tan\ang{12.0}$. Con $\tan\ang{12.0} =
0.21256$: la base queda por debajo de
\qty{200}{m} cuando $u > 80.16/0.21256 = 377.1$, y el techo por encima cuando
$u < 260.16/0.21256 = 1224.0$.

\textbf{21.} $x = 3500 \pm 377.1$ y $3500 \pm 1224.0$: el túnel está en la caliza desde
\qty{2276}{m} hasta \qty{3123}{m}, y desde \qty{3877}{m} hasta \qty{4724}{m}.

\textbf{22.} A lo largo de una línea horizontal, la base y el techo de la caliza, separados
\qty{180}{m} en vertical y ambos inclinados con $\tan\beta$, se cruzan a una distancia horizontal
$h/\tan\beta$: $180/0.21256 = \qty{846.8}{m}$ en cada flanco, lo mismo, con el redondeo, que
$1224.0 - 377.1 = 846.9$.

\textbf{23.} $2\times846.8 = \qty{1693.7}{m}$: el túnel perfora
\textbf{unos \qty[text-series-to-math=true]{1694}{m} de caliza}, más de una quinta parte de sus
\qty{8}{km} de longitud, en dos tramos de \qty{847}{m} a cada lado del núcleo de arenisca del
\omterm{def:b1:geological-maps:fold}{anticlinal}.
\end{solution}
